\clearpage\begingroup\SourceLeaf{412b}\ArithmeticHead{81}{R}{ELEMENTARE UND ANALYTISCHE GEOMETRIE.}
\renewcommand{\thefootnote}{\arabic{footnote})}\setcounter{footnote}{2}
\begin{center}
\textbf{27.}\par\smallskip
\textls[100]{Sphärische Trigonometrie.}
\end{center}
\noindent\hspace*{\parindent}\textsc{De Gua}\footnote{J. P. \textsc{de Gua}, \textit{Trigonométrie sphérique}, Mém. de l’Acad., année 1783, Paris 1786, S. 291.} und \textsc{Lagrange}\footnote{J. L. \textsc{Lagrange}, \textit{Solution de quelques problèmes relatifs aux triangles sphériques}, Journal de l’école polytechnique, cahier 6. 1798, S. 279, Oeuvres, t. 7, S. 329.} hatten gezeigt, dass die Kosinusformel zum
Aufbau der ganzen sphärischen Trigonometrie ausreicht; ihre Ableitungen
gelten indessen nur für Bogen, die nicht grösser als \(90^0\) sind. In dem \textit{Zusatz}
VII (1810) zu \textsc{Carnots} \textit{Geometrie der Stellung} (2. Teil, S. 373, W. IV, S. 401)
hat \textsc{Gauss} diese Lücke ausgefüllt und Bogen bis zu \(180^0\) zugelassen, wie sie
in der Praxis tatsächlich vorkommen\footnote{Für die rechtwinkligen Dreiecke hatte schon \textsc{Klügel} diese Erweiterung vorgenommen, \textit{Analytische Trigonometrie}, Braunschweig 1770.}. Aber schon in der \textit{Theoria motus corporum coelestium}, die 1809 erschienen war, hatte er im Art. 54 auf die allgemeinste
Auffassung des sphärischen Dreiecks hingewiesen, bei der weder Seiten
noch Winkel irgend welchen Beschränkungen unterworfen seien; die ausführlichere
Darstellung, die er in Aussicht stellte, ist aber weder veröffentlicht
\par\smallskip\noindent{\small X\textsubscript{2} Abh. 4.\hfill 11}\par
\clearpage\SourceLeaf{413}\ArithmeticHead{82}{L}{STÄCKEL, GAUSS ALS GEOMETER.}\setcounter{footnote}{0}
\noindent worden, noch hat sich im Nachlass etwas darüber gefunden. Später hat
\textsc{Möbius}, auch hier in den Spuren von \textsc{Gauss} wandelnd, die Untersuchung für
Bogen und Winkel durchgeführt, die bis \(360^0\) reichen\footnote{A. F. \textsc{Möbius}, \textit{Über eine neue Behandlungsweise der analytischen Sphärik}, Abhandlungen bei Begründung der Königl. Sächs. Gesellschaft der Wissenschaften, herausgegeben von der \textsc{Jablonowski}schen Gesellschaft d. W., Leipzig 1846, S. 45, Werke II, S. 1.}, zur vollen Allgemeinheit
ist aber erst \textsc{Study} (1893) gelangt\footnote{E. \textsc{Study}, \textit{Sphärische Trigonometrie, orthogonale Substitutionen und elliptische Funktionen}, Leipziger Abhandlungen, Bd. 21, 1893.}.
\par
Die vier Fundamentalformeln der sphärischen Trigonometrie lauten in
der Gestalt, die \textsc{Gauss} sich zu seinem Gebrauch aufgezeichnet hatte und die
er für die angemessenste hielt (Brief an \textsc{Schumacher} vom 26. September 1844,
Br. G.-Sch., IV, S. 310):
\[
\begin{gathered}
\cos a=\cos b\cos c+\sin b\sin c\cos A,\\
\sin a\sin B=\sin b\sin A,\\
\cos A\cos c=\operatorname{cotang}b\sin c-\operatorname{cotang}B\sin A,\\
\cos A=-\cos B\cos C+\sin B\sin C\cos a.
\end{gathered}
\]
\noindent Sie sind nebst den zugehörigen, ebenfalls von \textsc{Gauss} angegebenen Differentialformeln
in die Sammlung von Hülfstafeln aufgenommen worden, die \textsc{Warnstorff}
1845 als neue Ausgabe der von \textsc{Schumacher} 1822 veröffentlichten
Tafeln herausgegeben hat\footnote{G. H. L. \textsc{Warnstorff}, \textit{Sammlung von Hülfstafeln}, Altona 1845, S. 132. Die Formeln werden dort nicht ausdrücklich als von \textsc{Gauss} herrührend bezeichnet, während das bei den anderen Beiträgen von \textsc{Gauss} geschehen ist, z. B. bei den Tafeln für barometrisches Höhenmessen; vgl. W. IX, S. 456.}. Man findet hier auch eine Anweisung, die dritte
Formel dem Gedächtnis einzuprägen, die \textsc{Gauss}, wie \textsc{Wittstein} berichtet\footnote{Th. \textsc{Wittstein}, \textit{Lehrbuch der Elementar-Mathematik}, 2. Band, 2. Abteilung, Hannover 1862, S. 146: die betreffende Stelle ist abgedruckt W. X\textsubscript{1}, S. 457.},
seinen Zuhörern mitzuteilen pflegte.
\par
Im Art. 54 der \textit{Theoria motus} (1809) hatte \textsc{Gauss} ohne Beweis vier Gleichungen
zwischen den sechs Stücken eines sphärischen Dreiecks angegeben,
die er als nützlich für die Auflösung eines solchen Dreiecks bezeichnete, wenn
eine Seite und die anliegenden Winkel gegeben sind. Gefunden hatte er
diese Gleichungen, wie es scheint, auf dem Umwege von Betrachtungen über
die Frage, wie man die Gleichungen zwischen den Stücken eines sphärischen
Dreiecks auf Gleichungen zwischen den Stücken eines ebenen Dreiecks
\clearpage\SourceLeaf{414}\ArithmeticHead{83}{R}{ELEMENTARE UND ANALYTISCHE GEOMETRIE.}\setcounter{footnote}{0}
\noindent zurückführen könne (W. IV, S. 404). Später hat er in dem Brief an \textsc{Gerling} vom
18. Februar 1815 eine einfache Herleitung gegeben (W. VIII, S. 289); dabei
findet man zugleich die richtigen Vorzeichen der linken Seiten, die, wie \textsc{Gauss}
bereits in der \textit{Theoria motus} bemerkt hatte, bei Ausdehnung der Stücke über
\(180^0\) besonders bestimmt werden müssen.
\par
\textsc{Delambre} hat in der ausführlichen Besprechung der \textit{Theoria motus}, die
er in der \textit{Connaissance des temps pour l’an} 1812, Paris, juillet 1810, veröffentlicht
hat, darauf hingewiesen (S. 451), dass er jene Formeln schon im
Jahre 1807 bekannt gemacht habe\footnote{\textit{Connaissance des temps pour l’an} 1809, Paris avril 1807, S. 445. Auch \textsc{Delambre} hat die Formeln ohne Beweis mitgeteilt. Ein Beweis ist zuerst von K. B. \textsc{Mollweide} gegeben worden, der die Formeln selbständig gefunden hat, \textit{Zusätze zur ebenen und sphärischen Trigonometrie}, Monatliche Correspondenz, Bd. 18, 1808, S. 394.}. Er fügt hinzu: »Quand j’eus trouvé ces
formules, j’en cherchai des applications qui pouvaient être vraiment utiles;
n’en voyant aucune je les donnai simplement comme curieuses«, und wiederholt
dreimal, dass er ihnen die \textsc{Neper}schen Analogien vorziehe (S. 364, 370,
385). Eine Erfahrung von mehr als hundert Jahren hat gezeigt, dass die
»\textsc{Delambre}schen Gleichungen« für die Auflösung der sphärischen Dreiecke
wahrhaft nützlich sind\footnote{Vgl. E. \textsc{Hammer}, \textit{Lehr- und Handbuch der ebenen und sphärischen Trigonometrie}, 4. Aufl., Stuttgart 1916, S. 479, 481.}; im Besonderen werden sie in der Geodäsie bei der
Berechnung der \textsc{Soldner}schen (rechtwinklig-sphärischen) Koordinaten angewandt\footnote{Vgl. W. \textsc{Jordan}, \textit{Handbuch der Vermessungskunde}, Bd. III, 4. Aufl., Stuttgart 1896, S. 259.}.
\par
Für den \textsc{Legendre}schen Satz von der Zurückführung eines kleinen sphärischen
Dreiecks auf ein ebenes Dreieck mit eben so langen Seiten sei auf
den fünften Abschnitt dieses Aufsatzes (Nr. 30) verwiesen. Hier möge nur
noch die zierliche Lösung der Aufgabe erwähnt werden, den Ort der Spitze
eines sphärischen Dreiecks auf gegebener Grundseite und mit gegebenem
Inhalt zu finden, die \textsc{Gauss} in dem Briefe an \textsc{Schumacher} vom 6. Januar 1842
entwickelt hat (W. VIII, S. 293). Sie gehört in die Zeit der »geometrischen
Nachblüte«, aus der die Mehrzahl der Untersuchungen herrührt, über die in
diesem Abschnitt berichtet worden ist.
\begin{center}\rule{0.24\textwidth}{0.3pt}\end{center}
\par\smallskip\noindent\hfill{\small 11\textsuperscript{*}}\par
\clearpage\endgroup
